\section{Propagating the original source across quarter}\label{qi:source-proof}

We now choose the interval after the analytic barrier $\sigma$ in Lemma~\ref{qi:diagonal} has been fixed. Let $x_0=1/4$, choose a small $\delta>0$, and put
\[
 h_0=\delta/4,\qquad b_N=\lceil N(x_0+2\delta)\rceil,\qquad x_b=b_N/N.
\]
The letter $h_0$ is a width, distinct from the color index $h=Nt$.
Targets lie in $|r/N-x_0|\le h_0$. The whole interval from target to anchor remains in the regular matrix domain.

\subsection{A strict cumulative loss, uniform for complex colors}

At $t=0$ the relative logarithmic modulus is
\begin{equation}\label{qi:f}
 f(x)=\log|\mu_f(x,0)/\mu_s(x,0)|
          =2\log\cot(\pi x),\qquad
 f(x_0)=0,\quad f'(x)=-\frac{4\pi}{\sin(2\pi x)}.
\end{equation}
Shorten $\delta$ until $-8\pi\le f'\le-4\pi$ on the propagation interval. For $x\le x_0$, the possible positive contribution between $x$ and $x_0$ is at most $4\pi h_0^2$; the contribution from $x_0$ to $x_0+2\delta$ is at most $-8\pi\delta^2$. For $x\ge x_0$, integration directly from $x$ to $x_0+2\delta$ gives at most
$-2\pi(4\delta^2-h_0^2)$. Thus uniformly on the entire target interval,
\begin{equation}\label{qi:negative-action}
 \int_x^{x_b}\log|\mu_f(u,0)/\mu_s(u,0)|\dd u
      \le-7\pi\delta^2+O(N^{-1}).
\end{equation}
For every intermediate starting point $y$ between the lowest target and the anchor, the absolute value of this integral is at most
$C_0\delta^2+O(N^{-1})$, where $C_0=40\pi$ suffices.

The nonzero root ratio has an analytic logarithm on the fixed regular domain. Its even dependence on $t$ gives a uniform estimate
\begin{equation}\label{qi:color-continuity}
 \left|\Log(\mu_f/\mu_s)(x,t)
       -\Log(\mu_f/\mu_s)(x,0)\right|\le C_1|t|^2.
\end{equation}
Taking real parts controls the actual moduli for complex $t$, not just for real colors. Choose $\delta$ first with
$C_0\delta^2<\sigma/8$, and then choose $t_0>0$ with
$3C_1\delta t_0^2<\pi\delta^2$. Ordinary uniform Riemann estimates, nonvanishing roots, and $d_i=\mu_i+O(N^{-1})$ give
\begin{align}
 |G_f(r)/G_s(r)|&\le C e^{-6\pi\delta^2N}
             &&\text{for target }r,\label{qi:target-ratio}\\
 \left|\log|G_f(j)/G_s(j)|\right|&\le \sigma N/4+C
             &&\text{at every intermediate }j.\label{qi:partial-ratio}
\end{align}
The rounding of $b_N$ and all diagonal corrections contribute $O(1)$ to the logarithms. The two inequalities have different roles: \eqref{qi:target-ratio} suppresses the actual competing source, whereas \eqref{qi:partial-ratio} bounds conversions anywhere along its propagation.

\begin{figure}[ht]
 \centering\begin{tikzpicture}[x=1.55cm,y=1.7cm,>=Stealth,font=\small]
 \def\del{0.03125}
 \fill[orange!20]plot[domain=-.25:0,samples=30]
  (\x,{2*ln(cot(180*(.25+\del*\x)))})--(0,0)--(-.25,0)--cycle;
 \fill[blue!12]plot[domain=0:2,samples=40]
  (\x,{2*ln(cot(180*(.25+\del*\x)))})--(2,0)--(0,0)--cycle;
 \draw[->](-.7,0)--(2.65,0)node[right]{$x$};
 \draw[->](0,-1.2)--(0,.48)node[above]{$f(x)$};
 \draw[thick,blue!65!black]plot[domain=-.55:2.4,samples=80]
  (\x,{2*ln(cot(180*(.25+\del*\x)))});
 \draw[densely dashed](-.25,.26)--(-.25,-.26);
 \draw[densely dashed](.25,.26)--(.25,-.26);
 \draw[densely dashed](2,.15)--(2,-.9);
 \node[above left]at(-.25,.26){$x_0-h_0$};
 \node[above right]at(.25,.26){$x_0+h_0$};
 \node[above]at(2,.15){$x_b$};
 \node[below left]at(0,0){$x_0$};
 \node[blue!65!black]at(1.2,-.12){negative action};
 \draw[<-,thick](-.25,-1.08)--(2,-1.08)
   node[midway,below]{anchor-to-target propagation};
 \begin{scope}[shift={(4.1,-.18)},x=.8cm,y=.8cm]
   \draw[->](-1.15,0)--(1.25,0)node[right]{$\Re\mu$};
   \draw[->](0,-1.3)--(0,1.4)node[above]{$\Im\mu$};
   \draw[black!35](0,0)circle(1);
   \fill[blue!65!black](0,1)circle(2pt)
     node[right]{$\mu_s=\ii$};
   \fill[orange!70!black](0,-1)circle(2pt)
     node[right]{$\mu_f=-\ii$};
   \node[below,align=center]at(.15,-1.7){distinct roots\\at $x=x_0,\ t=0$};
 \end{scope}
\end{tikzpicture}

 \caption{The relative logarithmic modulus $f(x)=2\log\cot(\pi x)$ changes sign at $x_0=1/4$. Both roots are nonzero and distinct there. The short positive area below quarter is outweighed by the negative area up to the fixed anchor $x_b$ for every target in $[x_0-\delta/4,x_0+\delta/4]$. The proof additionally bounds the absolute action accumulated to every intermediate index, so a conversion at such an index is controlled. Horizontal distances in the drawing are rescaled; the fixed width is chosen only after the analytic remainder barrier.}\label{qi:action-figure}
\end{figure}

\subsection{Both initial components come from the finite sum}

Shrink $t_0$ further to lie in the color disk of the high-index result on a fixed interval containing the anchor and its next index. This preserves all preceding action inequalities. At $x_b$ the high-index result of \cite[Section 3.2--3.3]{Conway} applies on a fixed compact interval above quarter and uniformly for complex $|t|<t_0$. Set
\[
 X_r=(U_{Nt,r},U_{Nt,r+1})^{\mathsf T},\qquad
 \begin{pmatrix}a_{s,N}(t)\\a_{f,N}(t)\end{pmatrix}
       =V_N(x_b,t)^{-1}X_{b_N}.
\]
These are exact entries defined from \eqref{qi:source}. The high-index theorem and $V_N=V_0+O(N^{-1})$ give
\begin{equation}\label{qi:anchor}
 a_{s,N}(t)=
 \frac{C(x_b,t)}{\mu_s(x_b,t)}
       e^{NS(x_b,t)}(1+O(N^{-1})),\qquad
 |a_{f,N}/a_{s,N}|\le C/N.
\end{equation}
To check the projection, the high-index ratio is
$U_{Nt,b_N}/U_{Nt,b_N+1}=\mu_s(x_b,t)+O(N^{-1})$.
Therefore $X_{b_N}$ is
$U_{Nt,b_N}\mu_s(x_b,t)^{-1}(\mu_s(x_b,t),1)^{\mathsf T}$
plus a relative $O(N^{-1})$ vector. Applying $V_N^{-1}$ gives
\eqref{qi:anchor}. In particular $a_{s,N}$ is holomorphic and nonzero on the fixed smaller disk. No conclusion about source overlap was drawn from the eigenvalues alone, and the second component is not set to zero.

By \eqref{qi:interaction} and \eqref{qi:partial-ratio},
\[
 \|E_j\|\le C e^{-3\sigma N/4},\qquad
 G_r^{-1}Y_r=
 \left[I+O(Ne^{-3\sigma N/4})\right]
       \begin{pmatrix}a_{s,N}\\a_{f,N}\end{pmatrix}.
\]
Let $v_1(x,t)=[V_N(x,t)]_{1,s}$. The first component after restoring $V_NG_r$ is therefore
\begin{equation}\label{qi:selected-source}
 U_{Nt,r}=a_{s,N}(t)v_1(r/N,t)G_s(r)
                 (1+O(e^{-c_1N}))
\end{equation}
for target indices and some $c_1>0$.
Indeed the direct fast term has relative size at most
$C N^{-1}e^{-6\pi\delta^2N}$ by \eqref{qi:anchor} and \eqref{qi:target-ratio}. The interaction error was bounded with the full partial-action width; at the target its fast output is smaller again. Bounded endpoint frames and their nonzero selected first component preserve these estimates. This proves a relative assertion for the original finite source with both channels retained.

\subsection{A fixed complex interpolation of the selected product}

Let $l_N(x,t)=\Log d_s(x,t)$, continuing $\Log\mu_s$. On a smaller fixed tube these logarithms are uniformly bounded and satisfy
$|\partial_x^j l_N|\le B j! a^{-j}$. Put $K_N=\lfloor c_2N\rfloor$, with $c_2>0$ sufficiently small, and define
\begin{align}
 E_N^s(x,t)={}&N\int_x^{x_b}l_N(u,t)\dd u
       +\frac{l_N(x,t)-l_N(x_b,t)}2\notag\\
 &+\sum_{k=1}^{K_N}\frac{B_{2k}N^{1-2k}}{(2k)!}
 \{\partial_x^{2k-1}l_N(x_b,t)
                 -\partial_x^{2k-1}l_N(x,t)\}.
 \label{qi:EM}
\end{align}
Here the Bernoulli numbers use $z/(e^z-1)=\sum B_jz^j/j!$.
The Euler--Maclaurin formula with periodic Bernoulli remainder
\cite[equation (2.10.1)]{DLMF}, applied separately to real and imaginary parts, gives the remainder for the actual product from $r$ to $b_N-1$. The Fourier series of the periodic Bernoulli polynomial
\cite[equation (24.8.1)]{DLMF} bounds its absolute value by
$4(2k)!/(2\pi)^{2k}$. Combining this with the Cauchy derivative estimate bounds the remainder by
\[
 CN\left(\frac{K_N}{\pi aN}\right)^{2K_N}\le C'e^{-c_3N};
\]
this follows from the displayed Cauchy bound and
$|B_{2k}|\le4(2k)!/(2\pi)^{2k}$.
It holds uniformly for the physical index grid and complex colors. Hence the explicit function
\begin{equation}\label{qi:I}
 I_N(x,t)=a_{s,N}(t)v_1(x,t)e^{E_N^s(x,t)}
\end{equation}
is holomorphic and nonzero on a fixed smaller joint domain, and
\eqref{qi:selected-source} proves \eqref{qi:grid}.
The anchor is an explicitly specified finite coefficient, not an unknown closed invariant. There is no interpolation of a moving support cutoff.

\subsection{Recovering the same leading amplitude}

It remains to show that \eqref{qi:I} has the leading expression in
\eqref{qi:complex}, including its normalization. The selected diagonal column of \eqref{qi:diagonal-identity} obeys the actual scalar recurrence up to an exponentially small relative defect. Its one-step scalar ratio is
\[
 v_{s,N}(x,t)=d_s(x,t)\frac{v_1(x,t)}{v_1(x+e,t)}.
\]
Both rows of the exact matrix identity imply
\begin{equation}\label{qi:Riccati}
 v_{s,N}(x,t)
 =p(x,t,e)+\frac{q_c(x,t,e)}{v_{s,N}(x+e,t)}
                        +O(e^{-\sigma' N})
\end{equation}
on a fixed smaller tube, for some $\sigma'>0$.
For example the second row gives
$v_1(x+e)=d_s(x)[V_N(x)]_{2,s}+O(e^{-\sigma N})$;
substitution in the first row, at $x$ and $x+e$, gives
\eqref{qi:Riccati}. All divisions use uniformly nonzero factors.

We know $v_{s,N}=\mu_s+O(e)$ with derivative bounds, but have not assumed a formal power series for it. Write $v_{s,N}=\mu_s+w_N$.
Cauchy's estimate gives $w_N'=O(e)$, so expanding
\eqref{qi:Riccati} yields
\[
 (1+q_0/\mu_s^2)w_N
 =e\{p_1+q_1/\mu_s-q_0\mu_{s,x}/\mu_s^2\}+O(e^2).
\]
The coefficient is $1-\mu_f/\mu_s$, uniformly separated from zero and equal to $2$ at quarter and $t=0$. Consequently
\begin{equation}\label{qi:first-correction}
 v_{s,N}=\mu_s+e\nu+O(e^2).
\end{equation}
Since $v_1=\mu_s+O(e)$ and $v_1'=\mu_s'+O(e)$, this implies
\[
 d_s=\mu_s+e(\nu+\mu_s')+O(e^2).
\]
Insert this expansion directly into \eqref{qi:EM}. Its higher Bernoulli terms have total size $O(N^{-1})$ on the smaller complex tube, by the same factorial estimates used for the exponential remainder. The extra integral of $\mu_s'/\mu_s$ equals
$\Log\mu_s(x_b)-\Log\mu_s(x)$. Combining it with the half-endpoint term and $v_1(x)/v_1(x_b)$ gives
\begin{align}
 \frac{v_1(x,t)e^{E_N^s(x,t)}}{v_1(x_b,t)}
 ={}&\sqrt{\frac{\mu_s(x,t)}{\mu_s(x_b,t)}}\notag\\
 &\times\exp\left\{
 N\int_x^{x_b}\Log\mu_s\,\dd u
       +\int_x^{x_b}\frac{\nu}{\mu_s}\,\dd u\right\}
       (1+O(N^{-1})).\label{qi:selected-leading}
\end{align}
This calculation holds at complex $x$; it is not a uniqueness argument on a finite mesh. At the anchor,
$a_{s,N}v_1(x_b)=C(x_b,t)e^{NS(x_b,t)}(1+O(N^{-1}))$
by \eqref{qi:anchor}. Multiplication in \eqref{qi:selected-leading} proves
\eqref{qi:complex} with exactly the $C,S$ of \eqref{qi:SC}.
Normalized remainders are holomorphic and uniformly bounded, so Cauchy's formula supplies the derivative statement. This completes the proof of Theorem~\ref{qi:main}.

At $t=0$ the values inherited from the original high-index normalization are, on both sides of quarter,
\[
 S(x,0)=\pi\ii(1/2-x)^2,\qquad
 C(x,0)=\sqrt{\sin\pi x}\,e^{-3\pi\ii(1/2-x)/2}.
\]
No connection constant or parity-dependent factor has been chosen at the crossing.
